\documentclass[11pt,a4paper]{article}
\usepackage[utf8]{inputenc}
\usepackage[T1]{fontenc}
\usepackage[spanish]{babel}
\usepackage{hyperref}
\usepackage{geometry}
\geometry{margin=2.5cm}
\usepackage{parskip}

\title{003 el nombre real}
\author{Alejandro Valenzuela Rivera\\\small Universidad Autónoma de Nuevo León}
\date{\today}
\begin{document}
\maketitle

\begin{quote}
\textbf{Sesgo declarado.} Este trabajo se sostiene desde el sesgo declarado de su autor, y no finge neutralidad. Lo que afirma se puede tumbar, y ese es el punto: el falsificador está escrito, no escondido.
\end{quote}

\section*{Resumen}
El nombre se lee de un tiron; la expansion va adentro; el nombre que solo se deletrea es etiqueta.

\section{Por qué se sostiene}
Un nombre sin acto es una mentira dicha con buena diccion.

\section{De dónde viene}
La autoridad, y el repo: los nombres que ya tenían las piezas.

\section{El caso}
Ninguna: los nombres de este repo, que ya se leen de un tirón.

\section{Método}
Un nombre que dice el acto, y que se lee de un tiron.

\section{Límites, y qué lo tumba}
Un nombre que solo se puede usar deletreando su acronimo.

\section{Lo que queda abierto}
\textquestiondown{}Como se nombra lo que no se puede describir?

\section*{Cómo se leyó esta tesis}
Esto se escribió en el repositorio man-made-empathy, del que este candado es una pieza: candados/03 el nombre y la obra/003 el nombre real.md. La máquina que lo compuso lo hizo desde este mismo archivo, y no agregó nada que no estuviera escrito. Si algo de esto es falso, se dice en el mismo sitio, y se corrige con fecha.

\end{document}
